\section{Source Composition and Complete Reporting}
\label{sec:appendix_sources}

Table~\ref{tab:sources} identifies the frozen source strata and their native ALL-unit counts. The cited
artifacts identify upstream resources; the local inclusion and qualification rules remain
those of the frozen v5 construction. In particular, dataset-card counts and original-paper
sample counts must not be substituted for this table's evaluation denominators.

\begin{table}[htbp]
\centering
\small
\caption{Frozen source composition. Counts are non-reference native ALL units, totaling 6,104.}
\label{tab:sources}
\setlength{\tabcolsep}{5pt}
\begin{tabularx}{\linewidth}{@{}lXr@{}}
\toprule
ID & Source artifact & ALL units \\
\midrule
M01 & MedEinst \citep{chen2026medeinst} & 1,780 \\
M02 & MedPIC-Bench \citep{medpicbench} & 167 \\
M06 & HIV-treatment causal scenarios \citep{prosperi2026hiv} & 13 \\
M10 & CPV-MedQA \citep{benkirane2025bias} & 451 \\
M11 & Counterfactual Cultural Cues \citep{rezaei2026cultural} & 517 \\
M12 & FairMedQA / AMQA \citep{xiao2025amqa} & 584 \\
M14 & DiversityMedQA \citep{rawat2024diversitymedqa} & 454 \\
M17 & MedPerturb \citep{gourabathina2025medperturb} & 450 \\
M20 & MedCounterFact \citep{mo2026faithfulness} & 788 \\
M22 & MedRGB \citep{ngo2024medrgb} & 450 \\
M23 & BioRAB \citep{li2025biorab} & 450 \\
Total & Non-reference native units & 6,104 \\
\bottomrule
\end{tabularx}
\end{table}


The CPV citation describes the perturbation method; M10 uses the authors'
\href{https://huggingface.co/datasets/kenza-ily/medqa-cpv}{MedQA extension}, rather than
substituting the paper's primary question cohort. MedPIC-Bench is cited through its
available dataset card, with no inferred publication date. Version-specific reference
metadata and source-to-stratum mappings are recorded in the companion bibliography audit.

MedEinst derives its diagnostic material from DDXPlus \citep{ddxplus}; that upstream
provenance is retained alongside the MedEinst citation.

Complete machine-readable tables accompany the draft: seven selected baseline
configurations across ALL and R1--R5; all 21 original-to-terminal-head comparisons;
the catalog and constant-answer controls; and, for the later CF-MoA study, all six arms by configuration; all
11 sources; overlapping R1--R5 categories; all five input formats; all maintain/respond
pair scores; every recorded family interval; and attributed costs and coverage. The
source-level plot is a view of those complete tables. It does not define a new filtered
benchmark. Input-format scores use input denominators, whereas source and R-category
scores use native-unit denominators.
